\documentclass[12pt, twoside]{article}
%\documentclass[12pt, twoside]{article}
\usepackage[letterpaper, margin=1in, headsep=0.2in]{geometry}
\setlength{\headheight}{0.6in}
%\usepackage[english]{babel}
\usepackage[utf8]{inputenc}
\usepackage{microtype}
\usepackage{amsmath}
\usepackage{amssymb}
%\usepackage{amsfonts}
\usepackage[nomessages]{fp} %\FPeval{\var-name}{2*sin(pi/6)}
\usepackage{siunitx} %units in math. eg 20\milli\meter
\usepackage{yhmath} % for arcs, overparenth command
\usepackage{tikz} %graphics
\usetikzlibrary{quotes, angles, arrows, arrows.meta}
\usepackage{pgfplots}
\usepgfplotslibrary{fillbetween}
\pgfplotsset{compat=1.16}
\usepackage{graphicx} %consider setting \graphicspath{{images/}}
\usepackage{parskip} %no paragraph indent
\usepackage{enumitem}
\usepackage{multicol}
\usepackage{venndiagram}

\usepackage{fancyhdr}
\pagestyle{fancy}
\fancyhf{}
\renewcommand{\headrulewidth}{0pt} % disable the underline of the header
\raggedbottom
\hfuzz=2mm %suppresses overfull box warnings

\usepackage{hyperref}

\fancyhead[LO]{La Scuola d'Italia / Dr.\ Huson / IB Math 12 \\* 30 September 2026}
\fancyhead[RO]{\fbox{\begin{minipage}{6cm}\footnotesize Nickname:\\[0.5cm]\end{minipage}}}
\fancyhead[LE]{\fbox{\begin{minipage}{6cm}\footnotesize Nickname:\\[0.5cm]\end{minipage}}}
\fancyfoot[C]{\thepage}

\begin{document}
\subsubsection*{1.6 Classwork: Sampling bias and survey methodologies}

\emph{Answer in IB terminology. The sampling techniques are: simple random,
convenience, systematic, quota and stratified. Data are either discrete
or continuous.}

\begin{enumerate}[itemsep=0.15cm]

%--- Problem 1: IA survey plan; population; convenience bias; stratified allocation; systematic [8 marks] ---
%--- Source: Original (freshly authored, AA SL 4.1), modelled on EXM.2.SL.TZ0.1 and 21M.1.SL.TZ2.7(e) ---
\item[1.]

Alex is writing a mathematical exploration (IA) on the sleep habits of
the students at a high school. The number of students in each grade is
shown in the table.

\begin{center}
\setlength{\tabcolsep}{10pt}\renewcommand{\arraystretch}{1.2}
\begin{tabular}{|l|c|c|c|c|c|}
\hline
Grade & 9 & 10 & 11 & 12 & Total \\
\hline
Number of students & 84 & 78 & 68 & 70 & 300 \\
\hline
\end{tabular}
\end{center}

\begin{enumerate}[itemsep=0.1cm]
  \item Write down the population for Alex's investigation.
  \item Alex first plans to survey the $30$ students in the Grade 12
  common room at lunch. State the name of this sampling technique and
  explain why it is likely to produce a biased sample.
  \item Alex decides instead to take a stratified sample of $30$
  students by grade.
  \begin{enumerate}[label=(\roman*)]
    \item Show that $8$ students will be selected from Grade 10.
    \item Find the number of students selected from each of the
    other grades.
  \end{enumerate}
  \item To choose the Grade 12 students, Alex lists their names in
  alphabetical order, picks a random starting name among the first
  ten, and then selects every $10^{\text{th}}$ name on the list. State
  the name of this sampling technique.
\end{enumerate}

\begin{tikzpicture}
  \draw (-0.6,0) rectangle (14.6,10.6);
  \foreach \y in {1,2,...,10} \draw[dotted] (0,\y)--(14.1,\y);
\end{tikzpicture}

\newpage

%--- Problem 2: Discrete/continuous; quota; simple random; outlier as error; leading question [8 marks] ---
%--- Source: Original (freshly authored, AA SL 4.1) ---
\item[2.]

Sam is collecting survey data for an exploration on whether screen time
is related to hours of sleep. For each student Sam records the number of
apps on their phone and the time, in hours, they slept last night.

\begin{enumerate}[itemsep=0.1cm]
  \item State whether each variable is discrete or continuous.
  \begin{enumerate}[label=(\roman*)]
    \item The number of apps on a phone.
    \item The time slept last night.
  \end{enumerate}
  \item Sam stands at the school entrance and surveys the first $15$
  boys and the first $15$ girls who arrive. State the name of this
  sampling technique.
  \item Explain why the technique in part (b) is not the same as a
  stratified sample by gender.
  \item The school office gives Sam a numbered list of all $300$
  students. Describe how Sam could use this list to select a simple
  random sample of $30$ students.
  \item One student records a sleep time of $19$ hours. Suggest what
  Sam should do with this data value, and give a reason.
  \item One survey question reads: ``Don't you agree that using your
  phone late at night ruins your sleep?'' Explain why this question
  may produce biased data.
\end{enumerate}

\begin{tikzpicture}
  \draw (-0.6,0) rectangle (14.6,10.4);
  \foreach \y in {1,2,...,10} \draw[dotted] (0,\y)--(14.1,\y);
\end{tikzpicture}

\end{enumerate}
\end{document}
